
\subsection{Material, Section and Loading Inputs}
The seven ETABS models use M20 concrete, HYSD415/HYSD500 reinforcement, 125~mm floor slabs and the loading patterns presented below. Material properties, section dimensions and gravity and seismic actions are summarized in Tables~\ref{tab:material_properties}--\ref{tab:load_patterns}.

\begin{table}[!htbp]
\centering
\caption{Principal material properties used in the ETABS models.}
\label{tab:material_properties}
\small
\resizebox{\textwidth}{!}{%
\begin{tabular}{@{}lcccc@{}}
\toprule
Material & $f_c$/$f_y$ (MPa) & $E$ (MPa) & $\nu$ & Structural application \\
\midrule
M20 concrete & $f_c=20$ & 22360.68 & 0.20 & RC beams, columns and slabs \\
HYSD415 rebar & $f_y=415$ & 200000 & -- & Longitudinal/confining reinforcement definitions \\
HYSD500 rebar & $f_y=500$ & 200000 & -- & Longitudinal/confining reinforcement definitions \\
\bottomrule
\end{tabular}
}
\end{table}

\begin{table}[!htbp]
\centering
\caption{Principal section families used in the seven ETABS models.}
\label{tab:section_properties}
\small
\resizebox{\textwidth}{!}{%
\begin{tabular}{@{}lll@{}}
\toprule
Element & Section type & Nominal size \\
\midrule
Beam & M20 rectangular beam & 250$\times$355 mm and 250$\times$380 mm \\
Column & M20 rectangular column & 350$\times$350 mm; selected stair columns 400$\times$400 mm \\
Floor slab & M20 shell/slab & 125 mm \\
\bottomrule
\end{tabular}
}
\end{table}

\begin{table}[!htbp]
\centering
\caption{Load patterns and seismic cases used in the seven-model analyses.}
\label{tab:load_patterns}
\small
\begin{tabular}{@{}lp{8.5cm}@{}}
\toprule
Group & Model loading \\
\midrule
Gravity & Dead (self-weight), Floor Live, Roof Live, Staircase, Floor Finish \\
Wall & Peripheral Wall, Internal Wall, Parapet Wall \\
Seismic & EQx ULS, EQy ULS, EQx SLS, EQy SLS; 5\% accidental eccentricity was specified in the Regular model \\
Mass source & Dead + floor finish + wall loads at full factor, with 0.30 of Floor Live and Staircase in the Regular model \\
\bottomrule
\end{tabular}
\end{table}

\subsection{NBC 105:2025 Response Spectrum Basis}
The modal response-spectrum method uses the NBC 105:2025 spectral-shape-factor definition. Figure~\ref{fig:nbc105_mrsm_spectrum} reproduces the official code figure used as the spectral reference; the study demand cases use Soil Types C and D and the zone factors listed above.
\begin{figure}[!htbp]
\centering
\includegraphics[width=0.86\textwidth]{figures/ch1/nbc105_2025_mrsm_spectrum.png}
\caption{NBC 105:2025 response spectrum (Source: NBC 105:2025 \cite{nbc105_2025}).}
\label{fig:nbc105_mrsm_spectrum}
\end{figure}

The response-spectrum plot in Figure~\ref{fig:nbc105_mrsm_spectrum} shows the code spectral-shape basis used for modal response-spectrum analysis. The study applies this code-defined demand framework to the selected soil classes and seismic zone factors rather than deriving a separate project-specific response spectrum \cite{nbc105_2025}.

\FloatBarrier

\subsection{Modal Analysis and Mass-Participation Extraction}
Modal analysis was evaluated for the Regular, L, L1, L2, T, U and Non-Parallel building configurations. Twelve modes were considered for each seismic demand case. The main response measures are natural period, frequency, translational modal participation in the X and Y directions (UX and UY), and cumulative mass participation (SumUX and SumUY). These parameters indicate the relative contributions of the vibration modes to the building response.

The modal comparison in Chapter~4 uses Soil Type C at $Z=0.40$ to maintain a consistent demand case across the seven building configurations. The analyses also consider the remaining soil and seismic-zone combinations described in the Seismic Parameters section.

\begin{table}[!htbp]
\centering
\caption{Modal quantities extracted from the seven ETABS modal-participation workbooks.}
\label{tab:modal_extraction_fields}
\small
\begin{tabular}{@{}p{3.5cm}p{4.0cm}p{6.1cm}@{}}
\toprule
Quantity & Workbook field & Use in this study \\
\midrule
Modal period & Period & $T_1$, $T_2$, and $T_3$ are reported for the matched seven-model comparison. \\
Per-mode translational participation & UX, UY & Identifies the dominant X- and Y-translational modes. \\
Cumulative translational participation & SumUX, SumUY & Checks how many modes are required to exceed 90\% cumulative mass participation in each horizontal direction. \\
Frequency information & Frequency, CircFreq, Eigenvalue & Retained in the normalized electronic dataset for future verification. \\
\bottomrule
\end{tabular}
\end{table}

\FloatBarrier
